\section{Fase 3 - PatchAgent: loop di tool-use e toolkit}
\label{sec:fase3-toolkit}

\subsection{Il Loop di Tool-Use: \texttt{AzureClient.chat\_with\_tools}}
\label{subsec:3.4.3}
L'implementazione concreta e operativa del ciclo iterativo dell'agente
risiede all'interno del modulo \texttt{services/azure\_client.py}, ed è orchestrata
dal metodo \texttt{chat\_with\_tools()}. Questo componente è responsabile della gestione
a basso livello dell'interazione tra il \textit{Large Language Model}
e il toolkit di strumenti a sua disposizione.

\subsubsection{L'Architettura Function-Calling dell'API OpenAI}
\label{subsubsec:3.4.3.1}
Il fondamento è l'architettura \textit{Function Calling} dell'API \textit{Chat Completions} di OpenAI, un
protocollo per la comunicazione bidirezionale tra modello e codice applicativo. Quando l'esecuzione di uno
strumento è necessaria, l'API restituisce un oggetto \texttt{tool\_calls} anziché un messaggio di testo: una lista
di invocazioni, ciascuna con identificatore univoco (\texttt{tool\_call\_id}), nome della funzione
(\texttt{function.name}) e argomenti serializzati in JSON (\texttt{function.arguments}). Il codice applicativo
intercetta la richiesta, esegue lo strumento Python e inietta i risultati nella conversazione come messaggi con
ruolo \texttt{role: "tool"} e il corretto \texttt{tool\_call\_id}. Il metodo \texttt{chat\_with\_tools()} astrae
questo protocollo, come illustrato di seguito:

\begin{lstlisting}[language=Python, label={lst:chat_with_tools}]
def chat_with_tools(self, messages: list, tools: list, max_tool_rounds: int) -> tuple:
    """Gestisce il loop conversazionale con invocazione automatica dei tool."""
    for _ in range(max_tool_rounds):
        # Chiamata all'API con retry esponenziale
        response = self._call_with_retry(messages, tools)
        response_message = response.choices.message
        
        # Mutazione in-place della storia conversazionale
        messages.append(response_message)
        
        # Se non ci sono tool calls, il modello ha generato il testo finale
        if not response_message.tool_calls:
            return response_message.content, messages
            
        # Esecuzione sequenziale delle funzioni richieste
        for tool_call in response_message.tool_calls:
            func_name = tool_call.function.name
            try:
                # Deserializzazione degli argomenti generati dal modello
                args = json.loads(tool_call.function.arguments)
            except json.JSONDecodeError:
                # Fallback di sicurezza in caso di allucinazione sintattica
                args = {} 
                
            # Esecuzione fisica dello strumento (dispatch)
            result = self._execute_tool(func_name, args)
            
            # Mutazione in-place: aggiunta del risultato
            messages.append({
                "role": "tool",
                "tool_call_id": tool_call.id,
                "content": str(result)
            })
            
    # Raggiungimento del limite massimo di round
    return None, messages
\end{lstlisting}

Due aspetti implementativi meritano attenzione:

\begin{description}
    \item[La mutazione \textit{in-place} della lista \texttt{messages}:] la lista è passata per riferimento e
    modificata con \texttt{.append()} a ogni iterazione, così che il chiamante
    (\texttt{PatchAgent.\_run\_single\_attempt()}) abbia sempre l'intera storia conversazionale sincronizzata, anche
    in caso di terminazione prematura. L'oggetto \texttt{updated\_messages} restituito non è una copia, ma la stessa
    istanza mutata.

    \item[Il blocco try/except per \texttt{json.loads()}:] gli argomenti generati dal modello vanno deserializzati
    da JSON, ma l'LLM può produrre JSON malformato. Il blocco \texttt{try/except json.JSONDecodeError} intercetta
    l'anomalia e applica un fallback a dizionario vuoto (\texttt{args = \{\}}): lo strumento fallirà la validazione
    interna restituendo un errore testuale che, nel round successivo, l'LLM userà per correggere autonomamente
    l'invocazione. Soluzione architetturalmente superiore al crash.
\end{description}

\subsubsection{Il Budget di Tool-Call e la Prevenzione di Loop Infiniti}
\label{subsubsec:3.4.3.2}
Un agente libero di invocare strumenti in un ciclo \texttt{while} rischia di innescare loop infiniti, esaurendo il
budget dell'API senza convergere. Il parametro \texttt{max\_tool\_rounds} agisce da \textit{safety limit}, derivato
come \texttt{config.MAX\_AGENT\_STEPS * 3}. Con \texttt{MAX\_AGENT\_STEPS = 3}, l'agente dispone di 9 round. Il
moltiplicatore riflette l'anatomia della FSM: ciascuna macro-fase (\texttt{EXPLORE}, \texttt{PATCH}, \texttt{COMPILE},
\texttt{TEST}) richiede circa un round, ma la fase \texttt{COMPILE} può innescare cicli di correzione inline.
Moltiplicare il limite base fornisce un margine per assorbire un intero ciclo FSM con un paio di riparazioni
sintattiche al volo, mantenendo un tetto rigido contro i loop viziosi.

\subsubsection{Il Retry Esponenziale per Errori Transitori}
\label{subsubsec:3.4.3.3}
L'affidabilità di un sistema cloud dipende dalla resilienza agli errori di rete e al rate-limiting imposto dalle
API di Azure OpenAI. L'interazione è delegata al wrapper \texttt{\_call\_with\_retry()}, che implementa un backoff
esponenziale:

\begin{lstlisting}[language=Python, label={lst:call_with_retry}]
def _call_with_retry(self, messages: list, tools: list):
    """Esegue la chiamata API gestendo gli errori transitori di rate limit."""
    delays = [2, 4, 8]  # Backoff esponenziale in secondi
    
    for attempt_idx, delay in enumerate(delays):
        try:
            return self.client.chat.completions.create(
                model=self.model,
                messages=messages,
                tools=tools
            )
        except openai.RateLimitError as e:
            if attempt_idx == len(delays) - 1:
                # Esauriti i retry, propaga l'eccezione
                raise e
            # Attende prima di ritentare
            time.sleep(delay)
\end{lstlisting}

L'intervallo di attesa crescente (2, 4, 8 secondi) è il protocollo raccomandato per \texttt{RateLimitError}, ed è
vitale in modalità batch (es. \texttt{run\_all\_bugs.py}). Quando decine di thread paralleli vengono respinti per
rate-limit, un'attesa fissa li farebbe ripresentare in blocco contro i server (\textit{thundering herd}), causando
un nuovo collasso. Il backoff esponenziale scagliona le ritrasmissioni nel tempo, disperdendo i picchi di carico
e riducendo la probabilità di reimpattare contro i limiti.

\subsection{Il Toolkit dell'Agente: Sei Strumenti per la Riparazione Autonoma}
\label{subsec:3.4.4}
Il modulo \texttt{patch\_agent/tools.py} definisce l'arsenale operativo del \texttt{PatchAgent}: sei strumenti
esposti al modello tramite gli schemi JSON del \textit{Function Calling} e implementati come funzioni Python. Lo
stato è gestito con il pattern \textit{Context Object}: anziché passare singolarmente percorsi e classi di test a
ogni funzione, il contesto del tentativo corrente è aggregato in un oggetto mutabile condiviso:

\newpage
\begin{lstlisting}[language=Python, label={lst:tool_context}]
@dataclass
class ToolContext:
    bug_dir: Path
    patched_dir: Path
    test_classes: list[str]
\end{lstlisting}

L'iniezione del \texttt{ToolContext} riduce l'accoppiamento tra il dispatcher (\texttt{create\_tool\_executor()})
e le singole implementazioni, rendendo il toolkit estensibile: aggiungere strumenti non richiede modifiche alla
firma di routing.

\subsubsection{\texttt{read\_file}: Lettura del Codice Sorgente con Numerazione di Riga}
\label{subsubsec:3.4.4.1}
Il primo strumento è \texttt{read\_file}, per ispezionare i file del progetto durante la fase esplorativa.

\begin{lstlisting}[language=Python, label={lst:tool_read_file}]
def tool_read_file(context: ToolContext, file_path: str) -> str:
    """Legge il contenuto di un file sorgente applicando numerazione di riga."""
    target_path = context.patched_dir / file_path
    
    with open(target_path, 'r', encoding='utf-8') as f:
        content = f.read()
        
    # Troncamento per preservare il budget di token
    if len(content) > 30000:
        content = content[:30000] + "\n... [TRUNCATED]"
        
    lines = content.split('\n')
    # Numerazione destra-allineata a 4 cifre
    numbered_lines = [f"{i+1:4d} | {line}" for i, line in enumerate(lines)]
    return '\n'.join(numbered_lines)
\end{lstlisting}

Due dettagli sono cruciali. La formattazione \texttt{f"\{i+1:4d\} | \{line\}"} produce una numerazione di riga
allineata a 4 cifre: GPT-4, pre-addestrato su diff di GitHub con formattazione analoga, riesce così a riferirsi a
una riga per numero e a fornirla come parametro \texttt{target\_line} al patching. Il limite di 30.000 caratteri
evita che file Java monolitici (come le classi di JFreeChart) saturino il \textit{token budget}, costringendo
l'agente a una navigazione più mirata.

\subsubsection{\texttt{search\_symbol}: Ricerca Testuale nel Source Tree}
\label{subsubsec:3.4.4.2}
Quando il file in esame non contiene le risposte necessarie, l'agente può esplorare l'intero albero dei sorgenti
tramite \texttt{search\_symbol}:

\begin{lstlisting}[language=Python, label={lst:tool_search_symbol}]
def tool_search_symbol(context: ToolContext, pattern: str) -> str:
    """Esegue una ricerca testuale ricorsiva nel progetto target."""
    import subprocess
    cmd = ["grep", "-rnI", "--include=*.java", pattern, str(context.patched_dir)]
    
    result = subprocess.run(cmd, capture_output=True, text=True)
    lines = result.stdout.split('\n')
    
    # Limite ai primi 20 risultati per preservare il contesto
    return '\n'.join(lines[:20])
\end{lstlisting}

Lo strumento è un wrapper attorno a \texttt{grep -rnI --include=*.java}, che esegue una ricerca
\textit{case-sensitive} ricorsiva sui soli file Java; il flag \texttt{-I} ignora i file binari (es. \texttt{.class})
che genererebbero output illeggibile. In fase \texttt{EXPLORE} serve a localizzare l'implementazione di un metodo
citato nello stack trace, trovare le assegnazioni di una variabile \texttt{null} o individuare implementazioni
multiple di un'interfaccia. L'output è troncato a 20 risultati.

\subsubsection{\texttt{get\_method\_signature}: Verifica delle API Reali}
\label{subsubsec:3.4.4.3}
Per risolvere le ambiguità del sistema di tipi di Java, l'agente dispone di \texttt{get\_method\_signature}, che
ispeziona dinamicamente le firme dei metodi:

\begin{lstlisting}[language=Python, label={lst:tool_get_method_signature}]
def tool_get_method_signature(context: ToolContext, class_path: str, method_name: str) -> str:
    """Estrae le firme reali dei metodi pubblici per risolvere ambiguità di overload."""
    # Riutilizzo del modulo di estrazione per coerenza
    from logic_fl.ingredient_forge import _extract_public_methods
    
    target_file = context.patched_dir / class_path
    all_methods = _extract_public_methods(target_file)
    
    # Filtraggio mirato sul nome del metodo
    relevant_methods = [m for m in all_methods if method_name in m]
    return '\n'.join(relevant_methods)
\end{lstlisting}

Lo strumento riutilizza la funzione \texttt{\_extract\_public\_methods()} di \texttt{ingredient\_forge.py} (Fase 2):
condividere l'estrattore garantisce che le API esplorate dinamicamente abbiano la stessa struttura testuale di
quelle fornite negli ``Ingredients''. Il parametro \texttt{method\_name} filtra i risultati per ridurre il rumore
delle classi con molti metodi.

\subsubsection{\texttt{apply\_ast\_patch}: L'Atto di Riparazione Strutturale}
\label{subsubsec:3.4.4.4}
Questo strumento sancisce la transizione dallo stato esplorativo allo stato
esecutivo (\texttt{PATCH}), materializzando la modifica al codice:

\begin{lstlisting}[language=Python, label={lst:tool_apply_ast_patch}]
def tool_apply_ast_patch(context: ToolContext, file_path: str, method_name: str, new_method_body: str, target_line: int = None) -> str:
    """Sostituisce strutturalmente il metodo target nel file sorgente."""
    from patch_agent.ast_applier import apply_and_write
    
    target_path = context.patched_dir / file_path
    
    # Delega la sostituzione all'infrastruttura Tree-sitter
    success, diff = apply_and_write(
        file_path=target_path, 
        method_name=method_name, 
        new_source=new_method_body,
        target_line=target_line
    )
    
    if success:
        return f"[SUCCESS] Patch applied successfully.\nDiff:\n{diff}"
    else:
        return f"[FAILED] Failed to apply patch."
\end{lstlisting}

La funzione delega il lavoro pesante a \texttt{apply\_and\_write()} del modulo \texttt{ast\_applier.py}, cuore della
manipolazione AST/CST via Tree-sitter. Il diff unificato restituito al modello ha funzione puramente diagnostica:
offre una conferma visiva di ciò che è stato alterato, così che l'agente possa accorgersi di eventuali rimozioni
accidentali di righe ancor prima di compilare.

\newpage
\subsubsection{\texttt{compile\_check} e \texttt{run\_tests}: Gli Oracoli di Validazione}
\label{subsubsec:3.4.4.5}
Gli ultimi due strumenti, \texttt{compile\_check} e \texttt{run\_tests}, fungono da oracolo sintattico e semantico,
guidando le fasi \texttt{COMPILE} e \texttt{TEST}. Per garantire isolamento e tolleranza ai guasti (timeout, loop,
dipendenze), entrambi delegano l'esecuzione al sandbox di \texttt{sandbox\_evaluator.py} anziché eseguire comandi
shell diretti. \texttt{compile\_check} si affida a \texttt{compile\_patched\_source()}, troncando l'output di
\texttt{javac} a 3000 caratteri per isolare la radice dell'errore primario. \texttt{run\_tests} innesca
\texttt{evaluate()}, che riesegue la suite JUnit e distilla un feedback biforcato in \texttt{PASS},
\texttt{COMPILE\_ERROR} (disallineamenti non intercettati prima) o \texttt{TEST\_FAIL} (con stack trace), fornendo
all'agente le coordinate per riconsiderare la logica della patch.

